\section{Protocol Architecture}

The three products share a common \emph{platform layer}: a unified order book, a solver network, a risk engine, the Harvest Vault, a hook composition system, and a zero-knowledge proof pipeline. Products are applications on top of that layer---Borrow consumes the order book and settlement pipeline; Prime Trade adds the Hyperliquid custody model; Earn exposes the vault. The platform layer is extensible: additional products can be built on the same infrastructure without modifying the core contracts.

\subsection{Order Book and RFQ Symmetry}

All loan activity flows through a single order book. Borrower requests and solver quotes are structurally identical objects---both specify collateral, principal, rate, term, and execution path. When a borrower submits a request, it enters the book as a standard entry; when a solver responds, the response is an entry on the other side of the same book. The matching engine treats both symmetrically---there is no special-purpose RFQ subsystem separate from the order book.

This symmetry means that a request-for-quote is simply a way to populate the order book. When the book is liquid, a borrower can match against existing offers immediately. When the book is thin---as it will be at launch, or on niche markets---the borrower's request triggers solver quotes that flow into the book and are matched against the request through the same pipeline. From the protocol's perspective, the two cases are indistinguishable: the same objects, the same matching logic, the same settlement path. This gives borrowers a consistent experience regardless of market depth, and ensures that all liquidity---whether pre-posted or solicited---settles through one verifiable pipeline.

The order book is rate-type-agnostic: it expresses rate type, term, amount, and collateral as explicit per-entry terms, and can accommodate both fixed-rate and dynamic-rate configurations. The matching layer does not enforce a rate model---it settles whatever terms the solver and borrower agreed to. Shipped products in the current release are fixed-rate, fixed-term.

\subsection{Criteria-Based Matching}

The matching engine evaluates counterparty criteria alongside economic terms---a valid match requires both price compatibility and criteria satisfaction. Criteria are a typed set of attributes attached to an intent: counterparty allowlists, jurisdictional restrictions, KYC-tier requirements, or sanctions-screen attestations. The matcher treats a pair of entries as compatible only if each side's criteria are satisfied by the other side's attributes. There is no separate permissioned market and no permissioned deployment---there is one order book whose entries vary in how restrictive they are.

Criteria originate from three sources. \textbf{Self-asserted criteria} are declared by the entry's author at submission---for example, ``fill only against counterparties carrying attestation~$X$ from issuer~$Y$.'' \textbf{Counterparty attributes} are attached to a participant's identity object, sourced from on-chain attestations or signed claims by recognized issuers; the matcher reads these as part of intent evaluation. \textbf{Market-level criteria} are embedded in the market's program hashes---a market whose collateral-manager program requires every counterparty to carry an attestation from a specific issuer is itself a restricted market, but it shares the matching engine, risk engine, settlement pipeline, and vault infrastructure with unrestricted markets.

The matching state-transition program evaluates a criteria-compatibility predicate as part of the match validity check: given two entries $a$ and $b$, the predicate returns true if and only if every criterion in $a.\mathit{requires}$ is satisfied by $b.\mathit{attributes}$ and vice versa. This is a pure function of matching state, provable inside the existing zero-knowledge epoch program without an additional verifier. Failing the predicate produces a non-match---the entries are never paired, and no information about which participants carry which restrictions is revealed publicly.

Criteria can be public or committed-only. A public criterion (e.g.\ ``US-restricted'') is visible in the order book. A committed-only criterion is stored as a hash; the predicate evaluates against a witness inside the proof, revealing nothing beyond the boolean outcome. Counterparty attributes are typically committed-only---the matcher proves ``this counterparty satisfies the required attribute'' without exposing the full attribute set. This privacy posture is the primary knob for institutional adoption: regulated counterparties can participate in the unified order book without disclosing their compliance profile to every other participant.

An unrestricted intent---one with no criteria attached and no required counterparty attributes---matches against anyone who likewise carries no incompatible restrictions. The default is open. Criteria are opt-in constraints on either side. Adding criteria to the system does not change the experience of participants who do not use them; the unrestricted order book behaves exactly as described above.

\subsection{Markets}

\daloc{} is organized around isolated, immutable lending markets. Each market is identified by a deterministic hash of seven parameters:
\begin{equation}
  \texttt{marketId} = \texttt{keccak256}\!\left(c,\; \mathit{chainId}_p,\; p,\; h_{\mathit{rate}},\; h_{\mathit{liq}},\; h_{\mathit{di}},\; h_{\mathit{cm}}\right)
  \label{eq:marketid}
\end{equation}
where $c$ is the collateral asset, $\mathit{chainId}_p$ the principal chain, $p$ the principal asset, and $h_{\mathit{rate}}, h_{\mathit{liq}}, h_{\mathit{di}}, h_{\mathit{cm}}$ are zero-knowledge program hashes governing interest accrual, liquidation execution, debt-issuer state transition, and collateral-manager state transition respectively. Those parameters are fixed at creation and cannot be altered, giving participants a permanent guarantee of the rules governing their position; a different rate model, liquidation rule, or collateral policy constitutes a different market. Market-level criteria---when present---are encoded within the program hashes and therefore form part of the market identity: a market restricted to KYC-tier-2 counterparties has a different $h_{\mathit{cm}}$ and a different \texttt{marketId} from an otherwise identical unrestricted market, without requiring a separate registry or governance action. The protocol maintains no built-in interest model, no canonical oracle source, and no hardcoded liquidation logic at the contract level---all three live in the market's zero-knowledge programs. Markets are created permissionlessly through the hub contract. The hub is also the sole token rail: all user-facing ERC-20 transfers route through it, creating a single authorization point for capital flows, intent cancellations, and global pause.

Because market behavior is fully determined by its programs, different markets can encode fundamentally different collateral strategies. A \emph{classical} market locks collateral in the CDP for the loan's duration---the collateral serves solely as a liquidation hedge. A more capital-efficient market can define programs that permit managed rehypothecation: the collateral is deployed into yield-generating positions or used to back additional obligations, expanding the borrower's effective margin at the cost of additional risk. The zero-knowledge programs governing such markets are publicly visible---anyone can audit the logic---and cryptographically enforced, so positions always follow the rules the market defines, with no scope for deviation. The choice between conservative and aggressive collateral strategies is a market-definition decision, not a protocol-level toggle.

Deploying a market instantiates two coordinated contracts: a \emph{collateral manager} on the collateral chain and a \emph{debt issuer} on the principal chain. The collateral manager holds pledged assets, enforces the market's collateral policy (including any rehypothecation rules), and executes hook-composed pipelines for origination and unwind. The debt issuer mints and manages the loan's debt tokens, tracks interest accrual, and enforces repayment accounting. Both are deployed at deterministic addresses derived via \texttt{CREATE2} from the hub and the market identifier, so the same configuration yields the same contract addresses on every supported chain---one market identity without per-chain lookup tables. Individual loans are not created at market deployment; they are issued later, each with its own CDP and debt token, once settlement epochs commit them.

\subsection{Positions}

A position is a single loan. On the collateral side, each pledged asset creates an isolated \emph{CDP} (collateralized debt position)---one per loan, not shared across borrowers. There is no shared margin account. The CDP escrows whatever collateral or collateral-equivalent positions back the obligation; additional collateral can be topped up during the loan's life. On the funding side, the lender's claim is a \emph{debt token}: a transferable token representing the right to principal repayment plus accrued interest.

Isolation is a deliberate design choice that trades capital efficiency for risk containment. Unified-margin architectures aggregate a borrower's positions into one account and compute a single health factor across all venue exposure---collateral on the hub chain, equity at external trading venues, and PnL from open positions all feed one solvency metric. The efficiency gain from position netting is real, but the coupling is structural: if a venue fails, the global health factor drops and the entire account becomes liquidatable. \daloc{} takes the opposite approach: each loan's collateral sits in its own CDP, so one position's risk cannot cascade to another.

\subsection{Hook Composition}

Lending operations are expressed as atomic hook chains attached to settlement events. Each transfer in the settlement tree carries optional \emph{pre-hooks} and \emph{post-hooks}: the pre-hooks execute before the transfer (e.g.\ deploy a CDP), the transfer moves collateral or principal, and the post-hooks execute after (e.g.\ supply to an external venue, borrow, bridge). All steps must succeed or the entire operation reverts---there is no partial loan origination or partial unwind.

Hooks are the mechanism that makes \daloc{}'s multi-step pipelines possible within a single EVM transaction. A direct-path origination composes four hooks: escrow collateral in the CDP, supply collateral to the external lending venue, borrow principal, transfer principal to the borrower. A margin-path origination extends the chain: after borrowing, a bridge hook burns the principal via CCTP and a delivery hook deposits it on the destination chain. Each hook is a self-contained operation; the chain is assembled by the solver and committed into the settlement tree before on-chain execution.

This composability extends to collateral management. In classical markets, hooks simply lock collateral and supply it to a single venue. In markets whose zero-knowledge programs permit more aggressive strategies, hooks can compose arbitrarily complex pipelines: supply collateral to one venue, borrow an intermediate asset, supply that to a second venue, and route the proceeds---all within a single atomic transaction. The hook system imposes no limit on pipeline depth; the market's programs define what operations are authorized, and the settlement proof verifies that only authorized operations executed.

Borrowers sign high-level intents; solvers compete to attach the cheapest valid execution plan as a hook chain. Intent cancellation is always available---even when the protocol is paused---ensuring borrowers can revoke pending orders under any operational condition. Cancellation takes effect at the next epoch boundary; the settlement program reads on-chain cancellation state via storage proofs before matching.

\subsection{Epoch Settlement}

Time advances in \emph{epochs}: bounded settlement windows in which the off-chain engine matches intents, runs state-transition programs, and commits a snapshot to the chain. Each snapshot carries five Merkle roots: executable transfers, internal market state, newly finalized debt records, liquidation events, and authorized CDP deployments.\footnote{The internal market state uses a Poseidon-based Merkle Patricia Trie for ZK-friendly state management; the other four roots use keccak-hashed binary trees for on-chain proof verification.} Posting a new epoch requires a valid \emph{fact} in the facts registry---verified via zero-knowledge proof---attesting that the transition from the previous roots to the new ones was computed correctly by the market's registered programs.

Execution is inclusion-proof driven. Transfer intents, debt issuances, liquidations, and CDP deployments are leaves in Merkle trees; anyone can execute a committed leaf by supplying a proof against the epoch root stored on-chain. Epochs batch many loans and transfers into one verifiable step and ensure only consensus-approved state changes move funds.

\subsection{Solvers}

The Harvest Vault operator is the platform's canonical solver---it deploys vault capital into loans directly when idle capacity permits, or routes through external lending venues when vault liquidity is insufficient. Operating in solver mode, the vault operator participates in the order book like any other solver: responding to borrower requests with quotes, assembling hook chains, and committing them into settlement trees.

Independent solvers coexist on the same order book. A \emph{risk-taking} solver supplies the borrow from its own balance sheet, bearing principal risk in exchange for the spread against its own funding cost. A \emph{passive} solver routes through existing money-market liquidity---Aave~V3 in the initial release, with Morpho, Spark, and additional venues as expansion paths---by deploying the pledged collateral on established lending venues. Both modes produce structurally identical entries; a given request may receive responses from the vault operator, passive routers, and risk-taking principals, and the best terms win.

Solver compensation is spread-based. The core position contracts are asset-agnostic: they hold collateral and track obligations, but contain no references to specific lending venues, swap routers, or bridge protocols. Routing, pricing discovery, and capital allocation are solver-level concerns that live on top of the protocol rather than inside it. All external protocol interactions are composed from hook chains in the settlement commitment, making the platform extensible to new venues without modifying the core contracts.
